% GENERATED FILE. Edit canonical lessons, then run npm run generate:books.
% canonical-source-hash: fnv1a64:c89e9e02fcef5b18
% canonical-lessons: MR-C148-jad, MR-C148-mau, MR-C148-kadak, MR-C148-komat, MR-C148-suka

\chapter{Heavy, Soft, Hard, Lukewarm, Dry}
\label{ch:mr-heavy-soft-hard-lukewarm-dry}
\hlchaptermodality{148}

\begin{chapteropening}
\textbf{By the end of this chapter:} \emph{I can say heavy, soft, hard, lukewarm and dry.}

Complete the last lesson of chapter 148: I can say heavy, soft, hard, lukewarm and dry.
\end{chapteropening}

\section[\texorpdfstring{\mr{जड} (jaḍ)}{जड (jaḍ)}]{\mr{जड} (jaḍ) — heavy}

\label{lesson:MR-C148-jad}



\begin{quote}
Before the new one: say the Marathi for narrow, then the Marathi for thin.
\end{quote}

\subsection*{You'll want to know: \mr{जड}}
\textbf{\mr{जड}} — \emph{jaḍ} — \textquotedblleft{}heavy\textquotedblright{}.

\begin{cousinweb}
\textbf{\mr{जड}} (\emph{jaḍ}) is probably the word CDIAL lists as Marathi \emph{jaḍ}, \textquotedblleft{}cold, stupid\textquotedblright{}, under Sanskrit \emph{jaḍa}, \textquotedblleft{}stupid, dumb; cold\textquotedblright{}, a word it calls certainly of non-Aryan origin (CDIAL 5090). CDIAL does not give the sense \textquotedblleft{}heavy\textquotedblright{}.
\end{cousinweb}

\begin{grammarlens}[title={What you've built}]
One more word for this chapter.
\end{grammarlens}

\subsection*{Your turn}
\begin{itemize}
\raggedright
  \item \emph{Say it:} \emph{jaḍ}
  \item \emph{Say it:} \emph{jaḍ}, once more
  \item \emph{Say it:} say \emph{jaḍ}
  \item \emph{Recall:} say the Marathi for narrow, then the Marathi for thin, then say \emph{jaḍ} again
\end{itemize}

\begin{tcolorbox}[breakable,colback=teal!4,colframe=teal!35!black,title={Before you move on}]
What does \mr{जड} mean? (\textquotedblleft{}Heavy\textquotedblright{}.) Say it once more.
\end{tcolorbox}
\section[\texorpdfstring{\mr{मऊ} (maū)}{मऊ (maū)}]{\mr{मऊ} (maū) — soft}

\label{lesson:MR-C148-mau}



\begin{quote}
Before the new one: say the Marathi for thin, then the Marathi for heavy.
\end{quote}

\subsection*{You'll want to know: \mr{मऊ}}
\textbf{\mr{मऊ}} — \emph{maū} — \textquotedblleft{}soft\textquotedblright{}.

\begin{cousinweb}
\textbf{\mr{मऊ}} (\emph{maū}) is a tadbhava, a word inherited from Sanskrit and reshaped by sound change. It comes from \emph{mṛdu}, \textquotedblleft{}soft\textquotedblright{} (CDIAL 10292).
\end{cousinweb}

\begin{grammarlens}[title={What you've built}]
One more word for this chapter.
\end{grammarlens}

\subsection*{Your turn}
\begin{itemize}
\raggedright
  \item \emph{Say it:} \emph{maū}
  \item \emph{Say it:} \emph{maū}, once more
  \item \emph{Say it:} say \emph{maū}
  \item \emph{Recall:} say the Marathi for thin, then the Marathi for heavy, then say \emph{maū} again
\end{itemize}

\begin{tcolorbox}[breakable,colback=teal!4,colframe=teal!35!black,title={Before you move on}]
What does \mr{मऊ} mean? (\textquotedblleft{}Soft\textquotedblright{}.) Say it once more.
\end{tcolorbox}
\section[\texorpdfstring{\mr{कडक} (kaḍak)}{कडक (kaḍak)}]{\mr{कडक} (kaḍak) — hard, strict}

\label{lesson:MR-C148-kadak}



\begin{quote}
Before the new one: say the Marathi for heavy, then the Marathi for soft.
\end{quote}

\subsection*{You'll want to know: \mr{कडक}}
\textbf{\mr{कडक}} — \emph{kaḍak} — \textquotedblleft{}hard, strict\textquotedblright{}.

\begin{cousinweb}
The origin of \textbf{\mr{कडक}} (\emph{kaḍak}) is not settled. Neither CDIAL nor DEDR, searched by form and by meaning, has an entry for it.
\end{cousinweb}

\begin{grammarlens}[title={What you've built}]
One more word for this chapter.
\end{grammarlens}

\subsection*{Your turn}
\begin{itemize}
\raggedright
  \item \emph{Say it:} \emph{kaḍak}
  \item \emph{Say it:} \emph{kaḍak}, once more
  \item \emph{Say it:} say \emph{kaḍak}
  \item \emph{Recall:} say the Marathi for heavy, then the Marathi for soft, then say \emph{kaḍak} again
\end{itemize}

\begin{tcolorbox}[breakable,colback=teal!4,colframe=teal!35!black,title={Before you move on}]
What does \mr{कडक} mean? (\textquotedblleft{}Hard, strict\textquotedblright{}.) Say it once more.
\end{tcolorbox}
\section[\texorpdfstring{\mr{कोमट} (komaṭ)}{कोमट (komaṭ)}]{\mr{कोमट} (komaṭ) — lukewarm}

\label{lesson:MR-C148-komat}



\begin{quote}
Before the new one: say the Marathi for soft, then the Marathi for hard.
\end{quote}

\subsection*{You'll want to know: \mr{कोमट}}
\textbf{\mr{कोमट}} — \emph{komaṭ} — \textquotedblleft{}lukewarm\textquotedblright{}.

\begin{cousinweb}
\textbf{\mr{कोमट}} (\emph{komaṭ}) is a tadbhava, a word inherited and reshaped by sound change. It comes from \emph{*koṣma}, \textquotedblleft{}warm\textquotedblright{}, a form CDIAL reconstructs from its descendants (CDIAL 3552). CDIAL compares Sanskrit \emph{koṣṇa}, \textquotedblleft{}warm\textquotedblright{}.
\end{cousinweb}

\begin{grammarlens}[title={What you've built}]
One more word for this chapter.
\end{grammarlens}

\subsection*{Your turn}
\begin{itemize}
\raggedright
  \item \emph{Say it:} \emph{komaṭ}
  \item \emph{Say it:} \emph{komaṭ}, once more
  \item \emph{Say it:} say \emph{komaṭ}
  \item \emph{Recall:} say the Marathi for soft, then the Marathi for hard, then say \emph{komaṭ} again
\end{itemize}

\begin{tcolorbox}[breakable,colback=teal!4,colframe=teal!35!black,title={Before you move on}]
What does \mr{कोमट} mean? (\textquotedblleft{}Lukewarm\textquotedblright{}.) Say it once more.
\end{tcolorbox}
\section[\texorpdfstring{\mr{सुका} (sukā)}{सुका (sukā)}]{\mr{सुका} (sukā) — dry}

\label{lesson:MR-C148-suka}



\begin{quote}
Before the new one: say the Marathi for hard, then the Marathi for lukewarm.
\end{quote}

\subsection*{You'll want to know: \mr{सुका}}
\textbf{\mr{सुका}} — \emph{sukā} — \textquotedblleft{}dry\textquotedblright{}.

\begin{cousinweb}
\textbf{\mr{सुका}} (\emph{sukā}) is a tadbhava, a word inherited from Sanskrit and reshaped by sound change. It comes from \emph{śuṣka}, \textquotedblleft{}dried\textquotedblright{} (CDIAL 12548).
\end{cousinweb}

\begin{grammarlens}[title={What you've built}]
One more word for this chapter.
\end{grammarlens}

\subsection*{Your turn}
\begin{itemize}
\raggedright
  \item \emph{Say it:} \emph{sukā}
  \item \emph{Say it:} \emph{sukā}, once more
  \item \emph{Say it:} say \emph{sukā}
  \item \emph{Recall:} say the Marathi for hard, then the Marathi for lukewarm, then say \emph{sukā} again
\end{itemize}

\begin{tcolorbox}[breakable,colback=teal!4,colframe=teal!35!black,title={Before you move on}]
What does \mr{सुका} mean? (\textquotedblleft{}Dry\textquotedblright{}.) Say it once more.
\end{tcolorbox}
